% Lecture 9 notes: \input by notes.tex, \subimport'ed by the semester book.
% Exam questions (зачёт): 21 (линейные конгруэнтные последовательности, теорема о максимальном
% периоде), 22 (РСЛОС: рекуррентное уравнение, многочлен обратной связи, матричная форма).
\begin{lecture}{9}{Генераторы псевдослучайных чисел}

\section{Случайность в криптографии}

Почти каждый криптографический протокол расходует случайные числа: секретные показатели
в протоколе Диффи"--~Хеллмана, простые множители модуля RSA, одноразовые числа (nonce) в
подписях, ключи симметричных шифров, векторы инициализации. Стойкость протокола доказывается
в предположении, что эти числа выбраны \emph{равномерно и независимо}; если противник может
их угадать, никакая сложная задача теории чисел протокол уже не спасёт. Исторические примеры:
\begin{itemize}
  \item в браузере Netscape (1995) ключи SSL порождались из времени суток и номеров процессов "---
    их перебор занимал минуты;
  \item в OpenSSL из дистрибутива Debian (2006--2008) по ошибке осталась лишь энтропия номера
    процесса, и все ключи, созданные на таких системах, принадлежали множеству из нескольких
    десятков тысяч вариантов;
  \item в приставке Sony PlayStation~3 (2010) подписи ECDSA вычислялись с одним и тем же
    «случайным» числом $k$, что позволило найти секретный ключ подписи (см. семестр~2).
\end{itemize}

Источники истинной случайности (физический шум, моменты прерываний) медленны, поэтому на
практике короткое случайное \emph{зерно} растягивают детерминированным алгоритмом.

\begin{definition}
\emph{Генератор псевдослучайных чисел} (ГПСЧ) "--- это множество состояний $\mathcal{S}$,
функция перехода $F \colon \mathcal{S} \to \mathcal{S}$ и функция выхода
$G \colon \mathcal{S} \to \mathcal{Y}$. По начальному состоянию (зерну) $X_0$ генератор
вычисляет $X_{n+1} = F(X_n)$ и выдаёт $y_n = G(X_n)$, $n = 0, 1, 2, \dots$
\end{definition}

Требования к генератору зависят от применения.
\begin{itemize}
  \item \emph{Статистические}: выход должен проходить статистические тесты (равномерность,
    независимость пар и троек, распределение серий и т.\,п.), см.~\cite[разд.~3.3]{knuth-taocp2}.
    Этого достаточно для моделирования методом Монте"--~Карло, но не для криптографии.
  \item \emph{Криптографические} (\emph{тест следующего бита}): никакой полиномиальный алгоритм,
    получивший первые $i$ бит выхода, не предсказывает $(i+1)$-й бит с вероятностью, заметно
    большей $1/2$. По теореме Яо (1982, без доказательства) это равносильно тому, что выход
    неотличим от истинно случайной последовательности \emph{любым} полиномиальным
    статистическим тестом~\cite[разд.~5.5]{hac}.
\end{itemize}

Множество состояний конечно, поэтому любой ГПСЧ рано или поздно зацикливается.

\begin{proposition}\label{l9:rho}
Пусть $\abs{\mathcal{S}} = M$ и $X_{n+1} = F(X_n)$. Существуют $\mu \ge 0$ (\emph{предпериод}) и
$\lambda \ge 1$ (\emph{период}), $\mu + \lambda \le M$, такие, что $X_0, \dots, X_{\mu + \lambda - 1}$
попарно различны и $X_{n + \lambda} = X_n$ при всех $n \ge \mu$. Если $F$ "--- биекция, то $\mu = 0$
(последовательность \emph{чисто периодическая}).
\end{proposition}

\begin{proof}
Среди $M + 1$ элементов $X_0, \dots, X_M$ есть совпадающие; пусть $X_{\mu + \lambda}$ "--- первый
элемент, равный одному из предыдущих, $X_{\mu + \lambda} = X_\mu$. Тогда $X_0, \dots,
X_{\mu + \lambda - 1}$ различны, $\mu + \lambda \le M$, и по индукции
$X_{n + \lambda} = F^{n - \mu}(X_{\mu + \lambda}) = F^{n - \mu}(X_\mu) = X_n$ при $n \ge \mu$.
Если $F$ инъективна и $\mu > 0$, то из $F(X_{\mu - 1}) = X_\mu = X_{\mu + \lambda} =
F(X_{\mu + \lambda - 1})$ следует $X_{\mu - 1} = X_{\mu + \lambda - 1}$ "--- противоречие
с выбором $\mu$.
\end{proof}

\begin{center}
  \rhoshape[scale=0.6]{3}{8}
\end{center}
Траектория имеет форму буквы $\rho$ (та же картина появится в $\rho$-методе Полларда, семестр~2).
Период не может превышать числа состояний, и хороший генератор должен иметь период, недостижимый
на практике.

\checkyourself{Генератор хранит $64$-битное состояние. Каков максимально возможный период?
Можно ли за год «прокрутить» такой период, вычисляя $10^9$ значений в секунду?}

\section{Линейные конгруэнтные последовательности}

\begin{definition}
\emph{Линейная конгруэнтная последовательность} (ЛКП) с модулем $m \ge 2$, множителем $a$,
приращением $c$ ($0 \le a, c < m$) и зерном $X_0$ ($0 \le X_0 < m$):
\[
  X_{n+1} = (a X_n + c) \bmod m, \qquad n \ge 0.
\]
При $c = 0$ генератор называют \emph{мультипликативным}, при $c \ne 0$ "--- \emph{смешанным}.
\end{definition}

ЛКП "--- самый старый и самый простой ГПСЧ (Д.~Лемер, 1949). До сих пор это функция
\texttt{rand()} во многих стандартных библиотеках: например, в стандарте ANSI~C приведён
генератор с $a = 1103515245$, $c = 12345$, $m = 2^{31}$.

Обозначим $S_k(a) = 1 + a + \cdots + a^{k-1}$ (при $a \ne 1$ это $\frac{a^k - 1}{a - 1}$, при $a = 1$
это $k$).

\begin{proposition}\label{l9:lcg-closed}
$X_{n + k} \equiv a^k X_n + c\, S_k(a) \pmod m$ при всех $n, k \ge 0$.
\end{proposition}

\begin{proof}
Индукция по $k$: при $k = 0$ обе части равны $X_n$, и
$X_{n+k+1} \equiv a(a^k X_n + c S_k(a)) + c = a^{k+1} X_n + c (a S_k(a) + 1) = a^{k+1} X_n + c S_{k+1}(a)$.
\end{proof}

В частности, $X_k$ вычисляется за $\bigO(\log k)$ операций (бинарным возведением в степень, см.
лекцию~6) "--- так генератор «перематывают» вперёд.

Так как состояний $m$, период ЛКП не больше $m$. Поведение сильно зависит от параметров:
\[
\begin{array}{l|l|c}
  (a, c, m) & X_0 = 0,\ X_1, X_2, \dots & \text{период}\\ \hline
  (5, 3, 16) & 0, 3, 2, 13, 4, 7, 6, 1, 8, 11, 10, 5, 12, 15, 14, 9, 0, \dots & 16\\
  (3, 1, 16) & 0, 1, 4, 13, 8, 9, 12, 5, 0, \dots & 8\\
  (5, 2, 16) & 0, 2, 12, 14, 8, 10, 4, 6, 0, \dots & 8\\
  (7, 3, 10) & 0, 3, 4, 1, 0, \dots & 4\\
  (3, 1, 9) & 0, 1, 4, 4, 4, \dots & 1\ (\mu = 2)
\end{array}
\]

\begin{definition}
Говорят, что ЛКП имеет \emph{максимальный период}, если её период равен $m$.
\end{definition}

Если период равен $m$, то за период последовательность принимает все $m$ значений, поэтому
любое другое зерно $X_0'$ лежит на том же цикле, и последовательность с зерном $X_0'$ "--- сдвиг
исходной. Значит, свойство «период равен $m$» не зависит от зерна.

\subsection{Теорема о максимальном периоде}

\begin{theorem}[Халл"--~Добелл~\cite{hull-dobell}; Кнут, теорема~A]\label{l9:thmA}
ЛКП с параметрами $a, c, m$ имеет максимальный период $m$ тогда и только тогда, когда
\begin{enumerate}
  \item $\gcd(c, m) = 1$;
  \item $p \mid a - 1$ для каждого простого $p \mid m$;
  \item $4 \mid a - 1$, если $4 \mid m$.
\end{enumerate}
\end{theorem}

Для $m = 2^e$ условия означают: $c$ нечётно и $a \equiv 1 \pmod 4$. Генератор ANSI~C им
удовлетворяет: $12345$ нечётно, $1103515245 - 1 = 4 \cdot 275878811$. Генератор $(3, 1, 16)$ из
таблицы нарушает условие~3, $(5, 2, 16)$ "--- условие~1, $(7, 3, 10)$ и $(3, 1, 9)$ "---
условие~2.

Доказательство опирается на две леммы.

\begin{lemma}[неподвижная точка]\label{l9:fixed}
Если простое $p \mid m$ и $p \nmid a - 1$, то период меньше $m$.
\end{lemma}

\begin{proof}
По модулю $p$ число $1 - a$ обратимо; положим $x^* \equiv c (1 - a)^{-1} \pmod p$. Тогда
$a x^* + c \equiv x^* \pmod p$, т.\,е. если $X_n \equiv x^* \pmod p$, то и все последующие члены
сравнимы с $x^*$ по модулю $p$. Будь период равен $m$, последовательность прошла бы через все
вычеты, в том числе через вычет $\equiv x^* \pmod p$, после чего уже никогда не принимала бы
значений $\not\equiv x^* \pmod p$ "--- а на цикле, содержащем все вычеты, такие значения есть.
\end{proof}

\begin{lemma}\label{l9:lte}
Пусть $p$ простое, $a \equiv 1 \pmod p$, а при $p = 2$ даже $a \equiv 1 \pmod 4$. Тогда для всех
$n \ge 1$ показатель степени $p$ в $S_n(a)$ равен показателю $p$ в $n$: $v_p(S_n(a)) = v_p(n)$.
\end{lemma}

\begin{proof}
(i) Если $p \nmid n$, то $S_n(a) = \sum_{i<n} a^i \equiv n \not\equiv 0 \pmod p$.

(ii) Пусть $b \equiv 1 \pmod p$ (и $b \equiv 1 \pmod 4$ при $p = 2$). Покажем, что
$v_p(S_p(b)) = 1$. При $p = 2$: $S_2(b) = 1 + b \equiv 2 \pmod 4$. При нечётном $p$ запишем
$b = 1 + t$, $p \mid t$; по биному $(1 + t)^i \equiv 1 + it \pmod{t^2}$, поэтому
\[
  S_p(b) = \sum_{i=0}^{p-1} (1 + t)^i \equiv p + t \cdot \frac{p(p-1)}{2} \pmod{t^2}.
\]
Здесь $p^2 \mid t^2$ и $p^2 \mid t\,p\,\frac{p - 1}{2}$ ($\frac{p - 1}{2}$ "--- целое), значит,
$S_p(b) \equiv p \pmod{p^2}$.

(iii) Тождество $S_{pk}(a) = S_k(a)\, S_p(a^k)$ (разбейте сумму $\sum_{i < pk} a^i$ на $p$ блоков по
$k$ слагаемых) и (ii) для $b = a^k$ дают $v_p(S_{pk}(a)) = v_p(S_k(a)) + 1$. Индукцией по $v_p(n)$
с базой (i) получаем утверждение.
\end{proof}

\begin{proof}[Доказательство теоремы~\ref{l9:thmA}]
\emph{Необходимость.} (1) Пусть $d = \gcd(c, m) > 1$. При $X_0 = 0$ имеем
$X_{n+1} \equiv a X_n \pmod d$, т.\,е. все $X_n$ делятся на $d$ и принимают не более $m/d < m$
значений. (2)~Это лемма~\ref{l9:fixed}. (3)~Пусть $4 \mid m$; по (2) $a$ нечётно, и если
$a \equiv 3 \pmod 4$, то по модулю $4$
\[
  X_{n+2} \equiv a^2 X_n + c(a + 1) \equiv X_n \pmod 4,
\]
так как $a^2 \equiv 1$ и $a + 1 \equiv 0 \pmod 4$. Значит, вычеты $X_n \bmod 4$ принимают не более
двух значений, и последовательность не может пройти все вычеты по модулю $m$.

\emph{Достаточность.} По условию~2 каждое простое $p \mid m$ делит $a - 1$ и потому не делит $a$,
т.\,е. $\gcd(a, m) = 1$, отображение $x \mapsto ax + c$ "--- биекция $\Zn{m}$, и последовательность
чисто периодическая (утверждение~\ref{l9:rho}). Возьмём $X_0 = 0$; по
утверждению~\ref{l9:lcg-closed} $X_n \equiv c\, S_n(a)$, и так как $\gcd(c, m) = 1$,
\[
  X_n = 0 \iff m \mid S_n(a) \iff v_p(S_n(a)) \ge e_p \text{ для всех } p^{e_p} \,\|\, m.
\]
Для нечётных $p$, а также для $p = 2$ при $4 \mid m$ (тогда $a \equiv 1 \pmod 4$) по
лемме~\ref{l9:lte} $v_p(S_n(a)) = v_p(n)$. Если же $2 \,\|\, m$, то нужно лишь
$2 \mid S_n(a)$, а $S_n(a) \equiv n \pmod 2$ при нечётном $a$. В обоих случаях условие
равносильно $p^{e_p} \mid n$. Итак, $X_n = 0 \iff m \mid n$: последовательность впервые
возвращается в $0$ на шаге $m$, и её период равен $m$.
\end{proof}

\taskexample{lcg-full-period}

\begin{remark}
Максимальный период "--- необходимое, но далеко не достаточное условие качества. Генератор
$a = 1$, $c = 1$ имеет период $m$, но выдаёт $0, 1, 2, \dots$ При $m = 2^e$ и максимальном периоде
младший бит $X_n$ просто чередуется, а $k$ младших бит образуют ЛКП по модулю $2^k$ с периодом
$2^k$. Качество множителя проверяют \emph{спектральным тестом}~\cite[разд.~3.3.4]{knuth-taocp2}.
\end{remark}

\checkyourself{Докажите, что при $m = 2^e$ и максимальном периоде младшие биты $X_n \bmod 2$
чередуются: $0, 1, 0, 1, \dots$ или $1, 0, 1, 0, \dots$}

\subsection{Мультипликативные генераторы}

При $c = 0$ условие~1 теоремы~\ref{l9:thmA} нарушено ($0$ "--- неподвижная точка), и период меньше
$m$. Тем не менее такие генераторы популярны: они на одно сложение быстрее.

\begin{proposition}\label{l9:mult}
Пусть $c = 0$, $\gcd(a, m) = 1$, $X_0 \ne 0$ и $d = \gcd(X_0, m)$. Тогда последовательность чисто
периодическая с периодом $\ord_{m/d}(a)$ "--- порядком $a$ в группе $\Znx{m/d}$.
\end{proposition}

\begin{proof}
$X_n \equiv a^n X_0 \pmod m$, отображение $x \mapsto ax$ биективно. Далее
$X_n = X_0 \iff m \mid X_0 (a^n - 1) \iff \frac md \mid \frac{X_0}{d} (a^n - 1) \iff
\frac md \mid a^n - 1$, так как $\gcd(\frac{X_0}{d}, \frac md) = 1$.
\end{proof}

Наибольший порядок элемента группы $\Znx{m}$ обозначают $\lambda(m)$ (\emph{функция Кармайкла}).
Так как $\ord_{m/d}(a) \le \ord_m(a) \le \lambda(m)$, период мультипликативного генератора не
превосходит $\lambda(m)$ и равен $\lambda(m)$, если $\gcd(X_0, m) = 1$ и $\ord_m(a) = \lambda(m)$
(Кнут, теорема~B). Известно (без доказательства, см.~\cite[разд.~3.2.1.2]{knuth-taocp2} и
лекцию~7), что
\[
\begin{gathered}
  \lambda(2) = 1,\quad \lambda(4) = 2,\quad \lambda(2^e) = 2^{e-2}\ (e \ge 3),\quad
  \lambda(p^e) = p^{e-1}(p - 1)\ (p > 2),\\
  \lambda(p_1^{e_1} \cdots p_k^{e_k}) = \lcm\bigl(\lambda(p_1^{e_1}), \dots, \lambda(p_k^{e_k})\bigr).
\end{gathered}
\]
\begin{itemize}
  \item $m = p$ простое: наибольший период $p - 1$ достигается, когда $a$ "--- первообразный
    корень (лекция~7). Пример "--- генератор Парка"--~Миллера (MINSTD): $m = 2^{31} - 1$,
    $a = 7^5 = 16807$ "--- первообразный корень по модулю этого простого числа.
  \item $m = 2^e$: наибольший период $2^{e-2}$ (при $a \equiv 3, 5 \pmod 8$ и нечётном $X_0$).
    Печально известный генератор RANDU (IBM, 1960-е): $a = 65539 = 2^{16} + 3 \equiv 3 \pmod 8$,
    $m = 2^{31}$, период $2^{29}$. Однако $a^2 = 2^{32} + 6 \cdot 2^{16} + 9 \equiv 6a - 9
    \pmod{2^{31}}$, т.\,е. $X_{n+2} \equiv 6 X_{n+1} - 9 X_n$: все тройки последовательных
    значений лежат на 15 параллельных плоскостях в кубе.
\end{itemize}

\checkyourself{Почему при $m = 2^e$ мультипликативный генератор с чётным $X_0$ хуже, чем с
нечётным? Чему равен его период при $X_0 = 2^k X_0'$, $X_0'$ нечётно?}

\subsection{ЛКП предсказуемы}

Пусть противник видит несколько подряд идущих значений $X_0, X_1, X_2, \dots$ (например,
одноразовые числа протокола). Если модуль $m$ известен, то при обратимом $X_1 - X_0$
\[
  a \equiv (X_2 - X_1)(X_1 - X_0)^{-1}, \qquad c \equiv X_1 - a X_0 \pmod m,
\]
и все дальнейшие значения предсказываются. Неизвестный модуль тоже легко найти.

\begin{proposition}\label{l9:lcg-attack}
Пусть $t_n = X_{n+1} - X_n$ (разности целых чисел) и $u_n = t_{n+2} t_n - t_{n+1}^2$. Тогда
$t_{n+1} \equiv a t_n \pmod m$ и $m \mid u_n$ при всех $n$.
\end{proposition}

\begin{proof}
$t_{n+1} = X_{n+2} - X_{n+1} \equiv (a X_{n+1} + c) - (a X_n + c) = a t_n$. Поэтому
$t_{n+2} \equiv a^2 t_n$ и $u_n \equiv a^2 t_n^2 - (a t_n)^2 = 0 \pmod m$.
\end{proof}

Значит, $m$ делит $\gcd(u_0, u_1, \dots, u_k)$. Частные $u_n / m$ ведут себя как случайные
целые числа, а $k$ случайных чисел взаимно просты с вероятностью около $1/\zeta(k)$
(для $k = 3$ это $\approx 0{,}83$, для $k = 5$ "--- $\approx 0{,}96$). Поэтому уже по
$6$--$8$ значениям $\gcd(u_n)$ с высокой вероятностью равен $m$ (если нет "--- берём ещё
несколько значений). Найдя $m$, находим $a$ и $c$, как выше.

\taskexample{lcg-attack}

\begin{remark}
Не помогает и выдача только старших бит $X_n$: параметры такого «усечённого» генератора
восстанавливаются методами решёток (Фриз, Хостад, Каннан, Лагариас, Шамир, 1988; решётки "---
во втором семестре). Вывод: ЛКП хороши для моделирования, но \emph{никогда} не годятся для
криптографии.
\end{remark}

\checkyourself{Как восстановить $a$ и $c$, если $\gcd(X_1 - X_0, m) > 1$?}

\section{Регистры сдвига с линейной обратной связью}

Всюду в этом разделе арифметика ведётся в поле $\GF{2} = \{0, 1\}$: сложение "--- это XOR,
умножение "--- AND.

\begin{definition}
\emph{Регистр сдвига с линейной обратной связью} (РСЛОС, LFSR) длины $L$ с коэффициентами
$c_1, \dots, c_L \in \GF{2}$ порождает по начальному состоянию $(s_0, \dots, s_{L-1})$
последовательность
\begin{equation}\label{l9:rec}
  s_{n+L} = c_1 s_{n+L-1} + c_2 s_{n+L-2} + \cdots + c_L s_n, \qquad n \ge 0.
\end{equation}
\emph{Характеристический многочлен} (многочлен обратной связи) регистра
\[
  f(x) = x^L + c_1 x^{L-1} + \cdots + c_{L-1} x + c_L,
\]
\emph{многочлен связей} $C(x) = 1 + c_1 x + \cdots + c_L x^L = x^L f(1/x)$. Регистр
\emph{невырожден}, если $c_L = 1$, т.\,е. $f(0) = 1$.
\end{definition}

Аппаратно это $L$ ячеек памяти (битов). На каждом такте содержимое сдвигается на одну ячейку,
крайний бит $s_n$ выдаётся на выход, а в освободившуюся ячейку записывается сумма (XOR) битов
из ячеек с $c_j = 1$ ("--- \emph{отводов}). Такой регистр называют регистром в \emph{конфигурации
Фибоначчи}. Для $f(x) = x^4 + x + 1$ (т.\,е. $c_3 = c_4 = 1$, $s_{n+4} = s_{n+1} + s_n$):
\begin{center}
\begin{tikzpicture}[x=1.25cm,y=1cm,
    cell/.style={draw,thick,minimum width=1.1cm,minimum height=0.75cm},
    xorn/.style={draw,thick,circle,inner sep=0pt,minimum size=5mm}]
  \node[cell] (c3) at (0,0) {$s_{n+3}$};
  \node[cell] (c2) at (1,0) {$s_{n+2}$};
  \node[cell] (c1) at (2,0) {$s_{n+1}$};
  \node[cell] (c0) at (3,0) {$s_{n}$};
  \node[xorn] (x) at (2.5,1.2) {$+$};
  \draw[-{Stealth}] (c1.north) |- (x.west);
  \draw[-{Stealth}] (c0.north) |- (x.east);
  \draw[-{Stealth}] (x.north) -- ++(0,0.4) -| (-0.8,0) -- (c3.west);
  \draw[-{Stealth}] (c0.east) -- ++(0.8,0) node[right] {выход $s_n$};
  \node[left,font=\small] at (-0.8,0.8) {$s_{n+4}$};
\end{tikzpicture}
\end{center}
С начальным состоянием $(s_0, s_1, s_2, s_3) = (1, 0, 0, 0)$ получаем
\[
  1\,0\,0\,0\,1\,0\,0\,1\,1\,0\,1\,0\,1\,1\,1\ \ 1\,0\,0\,0\,\dots
\]
"--- последовательность с периодом $15 = 2^4 - 1$: регистр прошёл все ненулевые состояния.

\begin{remark}[конфигурация Галуа]
В программах удобнее другая схема: состояние "--- многочлен $r(x)$ степени $< L$, такт "---
$r(x) \gets x\, r(x) \bmod f(x)$ (сдвиг и, если вытолкнут старший бит, XOR с маской $f$), выход "---
старший коэффициент $u_n$ многочлена $x^n r_0(x) \bmod f$. Последовательность $u_n$ удовлетворяет
тому же соотношению~\eqref{l9:rec}: отображение $\varphi$, сопоставляющее вычету его коэффициент
при $x^{L-1}$, линейно, и
$u_{n+L} + c_1 u_{n+L-1} + \cdots + c_L u_n = \varphi\bigl(x^n r_0(x) f(x) \bmod f\bigr) = 0$.
Итак, схемы Фибоначчи и Галуа порождают одни и те же последовательности (при соответствующих
начальных состояниях).
\end{remark}

\subsection{Матричная форма и общий член}

Обозначим $S_n = (s_n, s_{n+1}, \dots, s_{n+L-1})^T \in \GF{2}^L$ "--- состояние регистра в момент
$n$. Соотношение~\eqref{l9:rec} в матричной форме:
\[
  S_{n+1} = A S_n, \qquad
  A = \begin{pmatrix}
    0 & 1 & 0 & \cdots & 0\\
    0 & 0 & 1 & \cdots & 0\\
    \vdots & & & \ddots & \vdots\\
    0 & 0 & 0 & \cdots & 1\\
    c_L & c_{L-1} & c_{L-2} & \cdots & c_1
  \end{pmatrix},
\]
где $A$ "--- \emph{сопровождающая матрица} многочлена $f$ (первые $L - 1$ строк сдвигают
состояние, последняя вычисляет $s_{n+L}$). Отсюда \emph{общий член}:
\begin{equation}\label{l9:general}
  S_n = A^n S_0, \qquad s_n = e_1^T A^n S_0,
\end{equation}
где $e_1 = (1, 0, \dots, 0)^T$. Матрица $A^n$ вычисляется бинарным алгоритмом (лекция~6) за
$\bigO(\log n)$ умножений матриц $L \times L$ над $\GF{2}$, т.\,е. $s_n$ находится за
$\bigO(L^3 \log n)$ битовых операций без вычисления предыдущих членов.

\begin{lemma}\label{l9:minpoly}
\begin{enumerate}
  \item $f(A) = 0$;
  \item $g(A) \ne 0$ для любого ненулевого многочлена $g$ степени $< L$;
  \item следовательно, $g(A) = 0 \iff f \mid g$ ($f$ "--- минимальный многочлен матрицы $A$);
  \item $A$ обратима $\iff c_L = 1$.
\end{enumerate}
\end{lemma}

\begin{proof}
(1) Первая координата вектора $A^k S$ "--- это $s_k$ для последовательности с начальным состоянием
$S$. Поэтому первая координата $f(A) S_n$ равна $s_{n+L} + c_1 s_{n+L-1} + \cdots + c_L s_n = 0$,
а $j$-я координата $f(A) S_0$ "--- это первая координата $A^{j-1} f(A) S_0 = f(A) S_{j-1}$, т.\,е.
тоже $0$. Так как $S_0$ произвольно, $f(A) = 0$.

(2) При $k < L$ строка $e_1^T A^k$ равна $e_{k+1}^T$ (матрица $A$ сдвигает состояние). Значит,
$e_1^T g(A) = \sum_{k < L} g_k e_{k+1}^T \ne 0$.

(3) Разделим $g$ на $f$ с остатком: $g = qf + r$, $\deg r < L$; по (1) $g(A) = r(A)$, и по (2)
$r(A) = 0 \iff r = 0$.

(4) При $c_L = 0$ первый столбец $A$ нулевой. При $c_L = 1$ рекуррентность обращается:
$s_n = s_{n+L} + c_1 s_{n+L-1} + \cdots + c_{L-1} s_{n+1}$, т.\,е. $S_n$ однозначно восстанавливается
по $S_{n+1}$.
\end{proof}

Из леммы получается и другая формула общего члена: если $x^n \bmod f = \sum_{k<L} r_k x^k$, то
$A^n = \sum_{k<L} r_k A^k$, и по~(2)
\[
  s_n = e_1^T A^n S_0 = \sum_{k=0}^{L-1} r_k s_k .
\]
Остаток $x^n \bmod f$ вычисляется бинарным алгоритмом за $\bigO(L^2 \log n)$ операций.

\taskexample{lfsr-period}

\subsection{Период}

Далее регистр невырожден ($c_L = 1$), а начальное состояние ненулевое (нулевое состояние
переходит в себя). Так как $f(0) = 1$, многочлен $x$ обратим по модулю $f$, и определён
\emph{порядок} $\ord_f(x)$ "--- наименьшее $e \ge 1$ с $x^e \equiv 1 \pmod f$, т.\,е. $f \mid x^e - 1$.

\begin{theorem}\label{l9:period}
Последовательность невырожденного РСЛОС чисто периодическая, её период $T$ не превосходит
$2^L - 1$ и делит $\ord_f(x)$.
\end{theorem}

\begin{proof}
По лемме~\ref{l9:minpoly}(4) $A$ "--- биекция $\GF{2}^L$, сохраняющая $0$, т.\,е. перестановка
$2^L - 1$ ненулевых векторов. По утверждению~\ref{l9:rho} последовательность $S_n$ чисто
периодическая с периодом $T \le 2^L - 1$. Пусть $e = \ord_f(x)$, $x^e - 1 = f h$; тогда
$A^e - I = f(A) h(A) = 0$, т.\,е. $S_e = S_0$, и минимальный период $T$ делит $e$.
\end{proof}

\begin{theorem}\label{l9:irred}
Если $f$ неприводим, то период любого ненулевого начального состояния равен $\ord_f(x)$, и
$\ord_f(x) \mid 2^L - 1$.
\end{theorem}

\begin{proof}
Кольцо $K = \GF{2}[x]/(f)$ "--- поле из $2^L$ элементов. Пусть $T$ "--- период, т.\,е.
$(A^T - I) S_0 = 0$. Положим $g = (x^T - 1) \bmod f$; по лемме~\ref{l9:minpoly} $g(A) = A^T - I$.
Если $g \ne 0$, то $g$ обратим в поле $K$: $gh \equiv 1 \pmod f$, откуда $g(A) h(A) = I$ и
матрица $A^T - I$ обратима, что противоречит $(A^T - I) S_0 = 0$, $S_0 \ne 0$. Значит,
$f \mid x^T - 1$, т.\,е. $\ord_f(x) \mid T$; вместе с теоремой~\ref{l9:period} $T = \ord_f(x)$.
Наконец, $x$ "--- элемент группы $K^*$ порядка $2^L - 1$, и по теореме Лагранжа его порядок делит
$2^L - 1$.
\end{proof}

\begin{definition}
Многочлен $f \in \GF{2}[x]$ степени $L$ называется \emph{примитивным}, если он неприводим и
$\ord_f(x) = 2^L - 1$ (т.\,е. $x$ "--- образующая циклической группы $K^*$).
\end{definition}

\begin{theorem}\label{l9:mseq}
РСЛОС длины $L$ порождает последовательность периода $2^L - 1$ (\emph{$M$-последовательность})
тогда и только тогда, когда его характеристический многочлен примитивен; в этом случае период
максимален для любого ненулевого начального состояния.
\end{theorem}

\begin{proof}
Если $f$ примитивен, то по теореме~\ref{l9:irred} период равен $\ord_f(x) = 2^L - 1$.
Обратно, пусть период $S_0$ равен $T = 2^L - 1$. Тогда $S_0, \dots, S_{T-1}$ "--- все ненулевые
векторы, и $A^T$ оставляет каждый из них на месте, т.\,е. $A^T = I$ и $f \mid x^T - 1$; по
теореме~\ref{l9:period} $\ord_f(x) = T = 2^L - 1$. Значит, вычеты $1, x, \dots, x^{2^L - 2}$
попарно различны, т.\,е. исчерпывают все $2^L - 1$ ненулевых вычетов по модулю $f$, и все они
обратимы. Если бы $f = gh$ с $0 < \deg g < L$, то вычет $g \ne 0$ был бы делителем нуля и не мог
быть обратим. Итак, $f$ неприводим, а значит, примитивен.
\end{proof}

\begin{proposition}[критерий примитивности]\label{l9:primcrit}
Неприводимый многочлен $f$ степени $L$ примитивен тогда и только тогда, когда
$x^{(2^L - 1)/q} \not\equiv 1 \pmod f$ для каждого простого $q \mid 2^L - 1$.
\end{proposition}

\begin{proof}
По теореме~\ref{l9:irred} $e = \ord_f(x)$ делит $2^L - 1$. Если $e \ne 2^L - 1$, то для некоторого
простого $q$ число $e$ делит $(2^L - 1)/q$, и $x^{(2^L - 1)/q} \equiv 1$. Обратное очевидно (как в
лекции~7 для первообразных корней).
\end{proof}

Неприводимость многочлена малой степени проверяют пробным делением на неприводимые многочлены
степени $\le L/2$: $x$, $x + 1$, $x^2 + x + 1$, $x^3 + x + 1$, $x^3 + x^2 + 1$, \dots{} Если
$2^L - 1$ "--- простое число Мерсенна ($L = 2, 3, 5, 7, 13, 17, 19, 31, \dots$), то любой
неприводимый многочлен степени $L$ примитивен. Число примитивных многочленов степени $L$ равно
$\varphi(2^L - 1)/L$ (без доказательства, см.~\cite[гл.~3]{lidl-niederreiter}):
\[
\begin{array}{c|cccccccc}
  L & 3 & 4 & 5 & 6 & 7 & 8 & 9 & 10\\ \hline
  2^L - 1 & 7 & 3 \cdot 5 & 31 & 3^2 \cdot 7 & 127 & 3 \cdot 5 \cdot 17 & 7 \cdot 73 & 3 \cdot 11 \cdot 31\\
  \text{примитивных} & 2 & 2 & 6 & 6 & 18 & 16 & 48 & 60
\end{array}
\]

\begin{example}
\begin{enumerate}
  \item $f = x^4 + x + 1$ неприводим (нет корней в $\GF{2}$ и не делится на $x^2 + x + 1$);
    $2^4 - 1 = 3 \cdot 5$, $x^3 \not\equiv 1$, $x^5 = x \cdot x^4 \equiv x^2 + x \not\equiv 1$:
    $f$ примитивен, период $15$ (см. выше).
  \item $f = x^4 + x^3 + x^2 + x + 1$ неприводим, но $x^5 - 1 = (x + 1) f$, т.\,е.
    $\ord_f(x) = 5$: \emph{каждое} ненулевое состояние имеет период $5$; например,
    $1000\,1\,1000\,1\dots$ "--- пространство из $15$ ненулевых состояний распадается на три цикла
    длины $5$.
  \item $f = x^4 + x^2 + 1 = (x^2 + x + 1)^2$ приводим, $\ord_f(x) = 6$; состояние $1000$ даёт
    период $6$ ($100010\,100010\dots$), а $1101$ "--- период $3$ ($110\,110\dots$).
\end{enumerate}
\end{example}

\begin{proposition}\label{l9:windows}
Если $f$ примитивен, то за период $M$-последовательности каждое ненулевое окно
$(s_n, \dots, s_{n+L-1})$ встречается ровно один раз. В частности, на периоде $2^{L-1}$ единиц и
$2^{L-1} - 1$ нулей.
\end{proposition}

\begin{proof}
Окна "--- это состояния $S_0, \dots, S_{2^L - 2}$, а они пробегают все ненулевые векторы ровно по
разу. Единиц столько, сколько ненулевых векторов с $s_n = 1$, т.\,е. $2^{L-1}$.
\end{proof}

$M$-последовательности обладают и другими «случайными» свойствами (постулаты Голомба:
распределение серий, двузначная автокорреляционная функция~\cite{golomb}), поэтому они широко
применяются в связи, навигации, тестировании схем.

\checkyourself{Докажите, что у вырожденного РСЛОС ($c_L = 0$) последовательность, вообще говоря,
не чисто периодическая. Как связаны её предпериод и период с разложением $f = x^k g$, $g(0) = 1$?}

\checkyourself{Почему для многочлена $x^4 + x^3 + x^2 + x + 1$ все ненулевые состояния имеют
\emph{один и тот же} период, а для $x^4 + x^2 + 1$ "--- разные?}

\section{Линейная сложность и восстановление РСЛОС}

Пусть РСЛОС используется как поточный шифр: $c_i = m_i \oplus s_i$. Противник, знающий кусок
открытого текста, знает и соответствующий кусок гаммы $s_i = c_i \oplus m_i$. Оказывается, $2L$
бит гаммы достаточно, чтобы восстановить регистр и всю дальнейшую гамму.

\begin{definition}
\emph{Линейная сложность} $\Lambda(s)$ (конечной или бесконечной) последовательности
$s$ "--- наименьшая длина РСЛОС (быть может, вырожденного), порождающего $s$. Последовательность
из нулей имеет сложность $0$.
\end{definition}

\begin{proposition}\label{l9:hankel}
Пусть бесконечная последовательность $s$ имеет линейную сложность $L$. Тогда векторы
$S_0, \dots, S_{L-1}$ линейно независимы, и коэффициенты $c_1, \dots, c_L$ однозначно находятся из
линейной системы над $\GF{2}$ по первым $2L$ битам:
\begin{equation}\label{l9:system}
  \begin{pmatrix}
    s_0 & s_1 & \cdots & s_{L-1}\\
    s_1 & s_2 & \cdots & s_L\\
    \vdots & & & \vdots\\
    s_{L-1} & s_L & \cdots & s_{2L-2}
  \end{pmatrix}
  \begin{pmatrix} c_L\\ c_{L-1}\\ \vdots\\ c_1 \end{pmatrix}
  =
  \begin{pmatrix} s_L\\ s_{L+1}\\ \vdots\\ s_{2L-1} \end{pmatrix}.
\end{equation}
\end{proposition}

\begin{proof}
Строки матрицы "--- векторы $S_0^T, \dots, S_{L-1}^T$, а уравнения "--- соотношение~\eqref{l9:rec}
при $n = 0, \dots, L - 1$. Пусть $\sum_{i<L} g_i S_i = 0$ для ненулевого $g = \sum g_i x^i$,
$\deg g = k < L$. Тогда $g(A) S_0 = 0$ и $g(A) S_n = A^n g(A) S_0 = 0$ при всех $n$; первая
координата даёт $\sum_{i \le k} g_i s_{n+i} = 0$, т.\,е. $s_{n+k} = \sum_{i<k} g_i s_{n+i}$ "---
последовательность порождается РСЛОС длины $k < L$, что противоречит определению $L$.
Значит, матрица невырождена, и решение единственно.
\end{proof}

\taskexample{lfsr-recover}

Если длина $L$ регистра известна лишь как верхняя граница, система~\eqref{l9:system} может
оказаться вырожденной (так бывает, когда $\Lambda(s) < L$). Тогда у неё несколько решений, но,
как показывает следующая лемма, все они порождают \emph{одну и ту же} последовательность.

\begin{lemma}[Месси~\cite{massey}]\label{l9:massey}
Если РСЛОС длины $L$ порождает $s_0, \dots, s_{N-1}$, но не $s_0, \dots, s_N$, то любой РСЛОС,
порождающий $s_0, \dots, s_N$, имеет длину $L' \ge N + 1 - L$.
\end{lemma}

\begin{proof}
Пусть $s_j = \sum_{i=1}^{L} c_i s_{j-i}$ при $L \le j \le N - 1$, но
$s_N \ne \sum_{i=1}^{L} c_i s_{N-i}$, и $s_j = \sum_{k=1}^{L'} d_k s_{j-k}$ при $L' \le j \le N$.
Допустим, $L' \le N - L$. Тогда при $1 \le i \le L$ индекс $N - i \ge L'$, и
\[
  \sum_{i=1}^{L} c_i s_{N-i} = \sum_{i=1}^{L} c_i \sum_{k=1}^{L'} d_k s_{N-i-k}
  = \sum_{k=1}^{L'} d_k \sum_{i=1}^{L} c_i s_{N-k-i} = \sum_{k=1}^{L'} d_k s_{N-k} = s_N,
\]
где в предпоследнем равенстве использовано, что $L \le N - k \le N - 1$ при $1 \le k \le L'$.
Противоречие.
\end{proof}

\begin{corollary}\label{l9:unique}
Если РСЛОС длин $L_1$ и $L_2$ порождают одни и те же первые $N \ge L_1 + L_2$ бит, то они
порождают одну и ту же бесконечную последовательность. В частности, если $s$ порождается
каким-то РСЛОС длины $\le L$, то любые два РСЛОС длины $\le L$, согласованные с первыми $2L$ битами
$s$, порождают $s$.
\end{corollary}

\begin{proof}
Пусть выходы $t$ и $t'$ впервые различаются в позиции $M \ge N$. Второй регистр порождает
$t_0, \dots, t_{M-1}$, но не $t_0, \dots, t_M$; по лемме~\ref{l9:massey} любой регистр,
порождающий $t_0, \dots, t_M$, имеет длину $\ge M + 1 - L_2 > L_1$ "--- а первый регистр длины
$L_1$ его порождает.
\end{proof}

Если $L$ неизвестно, можно решать системы~\eqref{l9:system} для $L = 1, 2, \dots$ ($\bigO(L^4)$
операций в сумме), но есть способ лучше.

\begin{algorithm}[Берлекэмп"--~Месси]
\begin{algorithmic}[1]
\Require биты $s_0, \dots, s_{n-1}$.
\Ensure линейная сложность $\ell$ и многочлен связей $C(x)$ кратчайшего РСЛОС.
\State $C(x) \gets 1$, $B(x) \gets 1$, $\ell \gets 0$, $m \gets -1$
\For{$N = 0, 1, \dots, n - 1$}
  \State $d \gets s_N + \sum_{i=1}^{\ell} c_i s_{N-i}$ \Comment{невязка}
  \If{$d = 1$}
    \State $T(x) \gets C(x)$, \quad $C(x) \gets C(x) + x^{N - m} B(x)$
    \If{$2\ell \le N$}
      \State $\ell \gets N + 1 - \ell$, \quad $m \gets N$, \quad $B(x) \gets T(x)$
    \EndIf
  \EndIf
\EndFor
\State \Return $(\ell, C(x))$
\end{algorithmic}
\end{algorithm}

\begin{theorem}[Месси, 1969; без доказательства]
После шага $N$ пара $(\ell, C)$ задаёт кратчайший РСЛОС, порождающий $s_0, \dots, s_N$.
Алгоритм выполняет $\bigO(n^2)$ битовых операций.
\end{theorem}

Идея: если текущий регистр $C$ ошибается на бите $s_N$ ($d = 1$), то регистр $B$, ошибавшийся
последний раз на шаге $m$, сдвинутый на $N - m$ позиций, ошибается \emph{ровно} на бите $s_N$, и
сумма $C + x^{N-m} B$ исправляет ошибку; лемма~\ref{l9:massey} показывает, что новая длина
$\max(\ell, N + 1 - \ell)$ не может быть меньше. Доказательство см.~\cite{massey},
\cite[разд.~6.2.3]{hac}. Характеристический многочлен найденного регистра $f(x) = x^\ell C(1/x)$.

\taskexample{berlekamp-massey}

\begin{remark}
Таким образом, у $M$-последовательности длины $L = 128$ период $2^{128} - 1$ астрономический, а
линейная сложность всего $128$: $256$ бит выхода определяют всё остальное. Большой период и
хорошая статистика не дают непредсказуемости.
\end{remark}

\checkyourself{Каков выход алгоритма Берлекэмпа"--~Месси на последовательности
$0, 0, \dots, 0, 1$ ($n - 1$ нулей и единица)? Почему это согласуется с леммой Месси?}

\section{Как сделать генератор стойким}

\subsection{Нелинейные комбинации РСЛОС}

Линейность "--- источник слабости, поэтому выходы нескольких РСЛОС комбинируют нелинейной
функцией или нерегулярно тактируют. Пример "--- генератор Геффе с тремя регистрами длин
$L_1, L_2, L_3$:
\[
  z = x_1 x_2 + (1 + x_2) x_3 \qquad (\text{если } x_2 = 1, \text{ выдаём } x_1, \text{ иначе } x_3).
\]
При примитивных многочленах и попарно взаимно простых длинах его линейная сложность равна
$L_1 L_2 + L_2 L_3 + L_3$ (без доказательства) и велика, но
$\Pr[z = x_1] = \Pr[z = x_3] = 3/4$ (проверьте перебором восьми наборов). \emph{Корреляционная
атака}: перебираем только начальные состояния первого регистра и выбираем то, при котором
его выход совпадает с $z$ примерно в $75\,\%$ позиций (для неверного состояния "--- около
$50\,\%$); так же находим третий регистр, затем второй. Вместо $2^{L_1 + L_2 + L_3}$ вариантов
нужно $\approx 2^{L_1} + 2^{L_2} + 2^{L_3}$.

Более удачные конструкции (фильтрующие и сжимающие генераторы, шифр A5/1 стандарта GSM с тремя
регистрами длин $19$, $22$, $23$ и мажоритарным тактированием) также были в разной степени
атакованы. Современные поточные шифры (ChaCha20, AES в режиме счётчика) строятся иначе.

\subsection{Генератор Блюма"--~Блюма"--~Шуба}

Пусть $N = pq$, где $p \equiv q \equiv 3 \pmod 4$ "--- большие простые. Зерно "--- квадрат
$x_0 = s^2 \bmod N$ случайного $s$, взаимно простого с $N$;
\[
  x_{i+1} = x_i^2 \bmod N, \qquad b_i = x_i \bmod 2 .
\]

\begin{theorem}[Блюм, Блюм, Шуб~\cite{blum-blum-shub}; без доказательства]
Если трудно отличать квадратичные вычеты по модулю $N$ от невычетов с символом Якоби $1$, то
генератор BBS проходит тест следующего бита. (Позднее было показано, что достаточно трудности
разложения $N$ на множители.)
\end{theorem}

Пример с игрушечными параметрами: $N = 11 \cdot 19 = 209$, $s = 3$, $x_0 = 9$:
\[
\begin{aligned}
  x_i&:\ 9, 81, 82, 36, 42, 92, 104, 157, 196, 169, 137, 168, 9, \dots,\\
  b_i&:\ 1, 1, 0, 0, 0, 0, 0, 1, 0, 1, 1, 0, \dots
\end{aligned}
\]
(период $12$). В отличие от ЛКП, знание $N$ и всех предыдущих битов не позволяет предсказать
следующий, но цена "--- одно возведение в квадрат по модулю $N$ (сотни цифр) на один бит, поэтому
BBS "--- скорее теоретический эталон. На практике криптографические ГПСЧ (Hash\_DRBG, HMAC\_DRBG,
CTR\_DRBG стандарта NIST SP~800-90A, генератор ядра Linux на основе ChaCha20) строят из хеш-функций
и блочных шифров, периодически подмешивая энтропию.

\checkyourself{Почему в генераторе BBS требуется $p \equiv q \equiv 3 \pmod 4$? (Подсказка: каждый
квадратичный вычет по модулю $N$ имеет ровно один квадратный корень, который сам является
квадратичным вычетом, т.\,е. возведение в квадрат "--- перестановка множества квадратичных
вычетов.)}

\checkyourselfsection

\begin{furtherreading}
  \readitem[разд.~3.1, 3.2.1, 3.3.4]{knuth-taocp2}{линейные конгруэнтные последовательности,
    теорема~A о максимальном периоде, мультипликативный случай, спектральный тест}
  \readitem[гл.~5, разд.~6.2--6.3]{hac}{криптографические ГПСЧ, тест следующего бита, генератор
    BBS (разд.~5.5.2); РСЛОС, линейная сложность, алгоритм Берлекэмпа"--~Месси, нелинейные
    комбинации}
  \readitem[гл.~3, 8]{lidl-niederreiter}{линейные рекуррентные последовательности над конечными
    полями: период, минимальный многочлен}
  \readitem[гл.~2]{paar-pelzl}{поточные шифры и регистры сдвига на простых примерах}
  \readitem{massey}{оригинальная статья: алгоритм Берлекэмпа"--~Месси и лемма о длине}
  \readitem{hull-dobell}{оригинальная статья: теорема о максимальном периоде}
  \readitem{wiki-lfsr}{схемы Фибоначчи и Галуа, таблицы примитивных многочленов}
  \readitem{coursera-boneh}{неделя~1: поточные шифры, ГПСЧ, непредсказуемость}
\end{furtherreading}
\end{lecture}
